\section{Datasets and Preprocessing}

The experiment deliberately uses two datasets with different classification scenarios to test the models under different text domains and label structures. Yahoo Answers Topics represents community question answering: short user questions and answers are assigned to 10 broad topics. AG News Sports vs Business represents binary news classification, where articles are separated into sports and business. The first setting is a multi-class informal Q\&A task, while the second is a cleaner two-class news task.

\begin{table}[H]
\centering
\caption{Full processed dataset statistics used for the experimental report.}
\label{tab:dataset_statistics}
\resizebox{\linewidth}{!}{%
\begin{tabular}{llrrrrr}
\toprule
Dataset & Scenario & Processed Rows & Labels & Train/Val/Test & Vocabulary & Avg. Length \\
\midrule
Yahoo Answers Topics & Community Q\&A & 6,000 & 10 & 4,200/900/900 & 34,981 & 95.15 \\
AG News Sports vs Business & News binary classification & 59,952 & 2 & 41,966/8,993/8,993 & 51,655 & 54.89 \\
\bottomrule
\end{tabular}
}
\end{table}


All statistics in Table~\ref{tab:dataset_statistics} are computed from the full processed dataset, meaning the train, validation, and test splits combined. This avoids reporting train-only statistics as if they described the whole dataset. Yahoo Answers contains 6,000 processed documents, and AG News contains 59,952 processed documents after cleaning and duplicate handling.

The preprocessing procedure applies lowercasing, tokenization, stopword removal, empty-text filtering, duplicate handling where necessary, and label standardisation. Traditional models and Word2Vec use the processed text representation. DistilBERT uses its own tokenizer during fine-tuning, so the text is not converted into sparse features before entering the model.

\begin{table}[H]
\centering
\caption{Top frequent words in the processed datasets.}
\label{tab:top_frequent_words}
\begin{tabular}{r l r l r}
\toprule
Rank & Yahoo Answers Word & Count & AG News Word & Count \\
\midrule
1 & the & 23,015 & the & 138,433 \\
2 & to & 15,561 & to & 78,777 \\
3 & a & 13,367 & a & 70,643 \\
4 & and & 13,166 & in & 61,120 \\
5 & i & 10,727 & of & 55,842 \\
6 & of & 10,635 & and & 40,535 \\
7 & you & 9,885 & 39 & 39,451 \\
8 & is & 9,471 & on & 38,209 \\
9 & in & 8,078 & for & 36,390 \\
10 & n & 7,550 & s & 33,259 \\
\bottomrule
\end{tabular}
\end{table}


Table~\ref{tab:top_frequent_words} summarises frequent content words after excluding standard English stopwords, non-alphabetic tokens, and obvious corpus artefacts from the recorded word-frequency files. The Yahoo Answers terms still reflect conversational question-answering language, whereas the AG News terms are more directly associated with sports and business reporting. Word frequency is therefore useful as a descriptive corpus summary, while the classification models rely on TF-IDF weighting, dense embeddings, or contextual token representations for prediction.

\begin{figure}[H]
    \centering
    \includegraphics[width=\linewidth]{figures/label_distribution.png}
    \caption{Label distributions computed from the full processed datasets.}
    \label{fig:label_distribution}
\end{figure}

Figure~\ref{fig:label_distribution} shows that both datasets are balanced. In Yahoo Answers, each of the 10 topic labels has 600 documents, so macro-averaged metrics are meaningful for comparing performance across classes. In AG News, the sports and business classes are almost equally represented, with 29985 and 29967 documents respectively. Therefore, later performance differences are more likely to reflect task difficulty and model representation than severe label imbalance.
